\documentclass[10pt,letterpaper]{article}
\usepackage[utf8]{inputenc}
\usepackage[margin=0.6in]{geometry}
\usepackage{amsmath,amssymb,amsthm}
\usepackage{graphicx}
\usepackage{xcolor}
\usepackage{enumitem}
\usepackage{booktabs}
\usepackage{tabularx}
\usepackage{multicol}
\usepackage{tcolorbox}
\usepackage{fancyhdr}
\usepackage{titlesec}
\usepackage[hidelinks]{hyperref}
\usepackage{mdframed}

% Colors
\definecolor{criticalcolor}{RGB}{220,20,60}
\definecolor{examplecolor}{RGB}{0,100,0}
\definecolor{warningcolor}{RGB}{255,140,0}
\definecolor{headercolor}{RGB}{25,25,112}
\definecolor{formulacolor}{RGB}{0,51,102}

% Header and Footer
\pagestyle{fancy}
\fancyhf{}
\fancyhead[L]{\small\textbf{ECO310 Final Exam - THE ULTIMATE GUIDE}}
\fancyhead[R]{\small Page \thepage}
\fancyfoot[C]{\small Master these concepts, ace the exam!}

% Section formatting
\titleformat{\section}{\LARGE\bfseries\color{headercolor}}{\thesection}{1em}{}[\titlerule]
\titleformat{\subsection}{\Large\bfseries\color{headercolor}}{\thesubsection}{1em}{}
\titleformat{\subsubsection}{\large\bfseries\color{formulacolor}}{\thesubsubsection}{1em}{}

% Custom boxes
\newtcolorbox{criticalbox}[1][]{
  colback=criticalcolor!5,
  colframe=criticalcolor,
  fonttitle=\bfseries,
  title={\textcolor{white}{★ CRITICAL - MASTER THIS FIRST}},
  #1
}

\newtcolorbox{examplebox}[1][]{
  colback=examplecolor!5,
  colframe=examplecolor!80,
  fonttitle=\bfseries,
  title={WORKED EXAMPLE},
  #1
}

\newtcolorbox{formulabox}[1][]{
  colback=formulacolor!5,
  colframe=formulacolor,
  fonttitle=\bfseries,
  title={KEY FORMULAS},
  #1
}

\newtcolorbox{warningbox}[1][]{
  colback=warningcolor!5,
  colframe=warningcolor,
  fonttitle=\bfseries,
  title={⚠ COMMON MISTAKE - AVOID THIS!},
  #1
}

\newtcolorbox{strategybox}[1][]{
  colback=blue!5,
  colframe=blue!70,
  fonttitle=\bfseries,
  title={PROBLEM-SOLVING STRATEGY},
  #1
}

% Compact lists
\setlist{nosep, leftmargin=*}

\begin{document}

% Title Page
\begin{titlepage}
\centering
\vspace*{1cm}

{\Huge\bfseries\color{criticalcolor} THE ULTIMATE\par}
\vspace{0.3cm}
{\Huge\bfseries ECO310 FINAL EXAM\par}
\vspace{0.3cm}
{\Huge\bfseries STUDY GUIDE\par}

\vspace{1cm}
{\Large\textit{Everything You Need to Ace This Exam}\par}

\vspace{1.5cm}

\begin{tcolorbox}[colback=yellow!10,colframe=orange,width=0.9\textwidth,arc=3mm]
\centering
\large
\textbf{One Guide. All Topics. Step-by-Step.}\\[0.3cm]
Choice Under Uncertainty + Markets Theory\\
From Zero to Hero in 3-4 Days
\end{tcolorbox}

\vspace{1cm}

{\large\textbf{What's Inside:}\par}
\begin{itemize}[leftmargin=2cm]
  \item ✓ Beginner-friendly explanations with real-world examples
  \item ✓ Every formula you need (color-coded by importance)
  \item ✓ 20+ fully worked examples
  \item ✓ Pattern recognition guide (when to use what technique)
  \item ✓ Common mistakes and how to avoid them
  \item ✓ Step-by-step problem-solving strategies
\end{itemize}

\vfill

{\large Created: March 31, 2026\par}
{\large Post-Midterm Content: Lectures 7-14\par}

\end{titlepage}

\tableofcontents
\newpage

% Quick Start Guide
\section*{START HERE: Quick Navigation}
\addcontentsline{toc}{section}{START HERE: Quick Navigation}

\begin{tcolorbox}[colback=yellow!20,colframe=orange!80,arc=3mm]
\textbf{First Time Learning?} Read in order: Section 1 → 2 → 3 → 4

\textbf{Need Quick Review?} Go straight to ``Formula Sheets'' (Section 6)

\textbf{Doing Practice Problems?} Use ``Pattern Recognition Guide'' (Section 5)

\textbf{Last-Minute Cramming?} Focus on Sections 1, 3, and 6 only
\end{tcolorbox}

\subsection*{The Big 2 Topics (Guaranteed on Exam)}

\begin{enumerate}[leftmargin=*]
  \item \textbf{Choice Under Uncertainty} (Section 1) - 50\% of final exam
    \begin{itemize}
      \item Expected utility, risk aversion, insurance, portfolio choice
      \item Master the certainty equivalent and risk premium
    \end{itemize}

  \item \textbf{Markets} (Section 2-4) - 50\% of final exam
    \begin{itemize}
      \item Cost functions, profit maximization
      \item Perfect competition vs. monopoly
      \item General equilibrium and welfare economics
    \end{itemize}
\end{enumerate}

\subsection*{How to Use This Guide}

\textbf{Day 1:} Read Section 1 (Uncertainty), do all examples, try 3 practice problems

\textbf{Day 2:} Read Sections 2-3 (Firms \& Monopoly), do all examples, try 3 practice problems

\textbf{Day 3:} Read Section 4 (General Equilibrium), review all sections, do full practice exam

\textbf{Day 4:} Review Section 6 (formulas), do another practice exam, make cheat sheet

\newpage

\section{CHOICE UNDER UNCERTAINTY - The First Half}

\begin{criticalbox}
\textbf{Why This Matters:} Core microeconomics topic. Covers 50\% of final exam material.

Master expected utility, risk aversion, and the certainty equivalent - these appear in almost every problem!
\end{criticalbox}

\subsection{The Big Picture: Why Uncertainty Matters}

\subsubsection{Real-World Examples}

Until now, we've assumed consumers know exactly what will happen. But in reality:
\begin{itemize}
  \item Should you buy insurance? (Pay certain premium to avoid uncertain loss)
  \item How to invest your savings? (Risky stocks vs. safe bonds)
  \item Should you accept a job offer with uncertain bonus?
\end{itemize}

\textbf{The key question:} How do people make decisions when outcomes are uncertain?

\subsection{Lotteries and Expected Value}

\subsubsection{What is a Lottery?}

A \textbf{lottery} is a random outcome with different prizes, each occurring with some probability.

\textbf{Notation:} $L = (x_1, p_1; x_2, p_2; \ldots; x_n, p_n)$
\begin{itemize}
  \item $x_i$ = outcome $i$ (usually money)
  \item $p_i$ = probability of outcome $i$
  \item $\sum_{i=1}^n p_i = 1$ (probabilities sum to 1)
\end{itemize}

\begin{examplebox}
\textbf{Example:} A coin flip lottery:
\[
L = (\$100, 0.5; \$0, 0.5)
\]

Win \$100 with 50\% chance, win \$0 with 50\% chance.
\end{examplebox}

\subsubsection{Expected Value}

The \textbf{expected value} is the average outcome if you played the lottery many times.

\begin{formulabox}
\textbf{Expected Value:}
\[
\boxed{EV(L) = \sum_{i=1}^n p_i x_i}
\]

This is just a weighted average of outcomes!
\end{formulabox}

\begin{examplebox}
\textbf{Problem:} Find the expected value of $L = (\$100, 0.5; \$0, 0.5)$

\textbf{Solution:}
\[
EV(L) = 0.5 \times 100 + 0.5 \times 0 = \boxed{\$50}
\]

On average, you expect to win \$50.
\end{examplebox}

\subsection{Expected Utility Theory}

\begin{criticalbox}
\textbf{KEY INSIGHT:} People don't maximize expected value - they maximize expected utility!

This is the foundation of everything in this section.
\end{criticalbox}

\subsubsection{The Problem with Expected Value}

Consider two lotteries:
\begin{itemize}
  \item Lottery A: Win \$1 million with 100\% probability
  \item Lottery B: Win \$10 million with 10\% probability, \$0 with 90\% probability
\end{itemize}

\textbf{Expected values:}
\begin{align*}
EV(A) &= \$1,000,000 \\
EV(B) &= 0.1 \times 10,000,000 + 0.9 \times 0 = \$1,000,000
\end{align*}

Same expected value! But most people prefer A (guaranteed money) over B (risky gamble).

\textbf{Why?} Because utility is not linear in money!

\subsubsection{Utility Functions}

A \textbf{utility function} $u(x)$ assigns a ``happiness level'' to wealth level $x$.

\begin{formulabox}
\textbf{Expected Utility:}
\[
\boxed{EU(L) = \sum_{i=1}^n p_i \cdot u(x_i)}
\]

People choose the lottery with highest expected utility, not highest expected value!
\end{formulabox}

\subsubsection{The Von Neumann-Morgenstern (vNM) Axioms}

Expected utility theory relies on four axioms about rational preferences:

\begin{enumerate}
  \item \textbf{Completeness:} Can compare any two lotteries
  \item \textbf{Transitivity:} If $L_1 \succ L_2$ and $L_2 \succ L_3$, then $L_1 \succ L_3$
  \item \textbf{Continuity:} Smooth preferences (no jumps)
  \item \textbf{Independence:} If $L_1 \succ L_2$, then mixing both with a third lottery $L_3$ preserves the preference
\end{enumerate}

\textbf{Result:} If someone satisfies these axioms, their preferences can be represented by expected utility!

\subsection{Risk Aversion}

\begin{criticalbox}
\textbf{MOST IMPORTANT CONCEPT:} Risk aversion determines all behavior under uncertainty!

Master this and you'll ace half the exam.
\end{criticalbox}

\subsubsection{What is Risk Aversion?}

A person is \textbf{risk averse} if they prefer a certain amount over a lottery with the same expected value.

\begin{examplebox}
\textbf{Example:} Consider lottery $L = (\$100, 0.5; \$0, 0.5)$ with $EV = \$50$

A risk-averse person prefers \$50 for sure over the lottery.

A risk-loving person prefers the lottery over \$50 for sure.

A risk-neutral person is indifferent.
\end{examplebox}

\subsubsection{Mathematical Definition}

\begin{formulabox}[colback=yellow!20,colframe=orange!80]
\textbf{RISK AVERSION AND CONCAVITY:}

A person with utility $u(x)$ is:
\begin{itemize}
  \item \textbf{Risk averse} if $u$ is \textbf{concave} (curves down): $u''(x) < 0$
  \item \textbf{Risk neutral} if $u$ is \textbf{linear}: $u''(x) = 0$
  \item \textbf{Risk loving} if $u$ is \textbf{convex} (curves up): $u''(x) > 0$
\end{itemize}

\textbf{Key property of risk aversion:}
\[
\boxed{u(EV(L)) > EU(L)}
\]

Utility of expected value exceeds expected utility!
\end{formulabox}

\subsubsection{Why Concavity = Risk Aversion}

\textbf{Intuition:} Concave utility means diminishing marginal utility of wealth.
\begin{itemize}
  \item Going from \$0 to \$50 gives a lot of utility
  \item Going from \$50 to \$100 gives less additional utility
  \item So losing \$50 hurts more than gaining \$50 helps!
  \item Therefore, prefer \$50 for sure over 50-50 gamble of \$0 or \$100
\end{itemize}

\begin{examplebox}
\textbf{Example:} Let $u(x) = \sqrt{x}$ (concave!)

Lottery: $L = (100, 0.5; 0, 0.5)$, so $EV(L) = 50$

\textbf{Expected utility:}
\[
EU(L) = 0.5 \times \sqrt{100} + 0.5 \times \sqrt{0} = 0.5 \times 10 + 0 = 5
\]

\textbf{Utility of expected value:}
\[
u(EV(L)) = u(50) = \sqrt{50} \approx 7.07
\]

Since $7.07 > 5$, we have $u(EV) > EU$, confirming risk aversion!
\end{examplebox}

\subsection{Certainty Equivalent and Risk Premium}

\begin{criticalbox}
\textbf{EXAM FAVORITE:} These concepts appear in 80\% of uncertainty problems!
\end{criticalbox}

\subsubsection{Certainty Equivalent (CE)}

The \textbf{certainty equivalent} is the guaranteed amount that makes you indifferent to the lottery.

\begin{formulabox}
\textbf{Certainty Equivalent:}
\[
\boxed{u(CE) = EU(L)}
\]

Find CE by solving: $u(CE) = \sum_i p_i \cdot u(x_i)$
\end{formulabox}

\textbf{Interpretation:} CE is the smallest amount you'd accept to give up the lottery.

\begin{examplebox}
\textbf{Problem:} You have $u(x) = \sqrt{x}$ and face lottery $L = (100, 0.5; 0, 0.5)$. Find CE.

\textbf{Solution:}

\textbf{Step 1:} Calculate $EU(L)$
\[
EU(L) = 0.5 \times \sqrt{100} + 0.5 \times \sqrt{0} = 5
\]

\textbf{Step 2:} Solve $u(CE) = EU(L)$
\[
\sqrt{CE} = 5 \implies CE = 25
\]

\textbf{Answer:} $CE = \$25$

You'd accept \$25 for sure instead of the lottery. Even though the lottery has expected value \$50, you only value it at \$25 because you're risk averse!
\end{examplebox}

\subsubsection{Risk Premium (RP)}

The \textbf{risk premium} measures how much you'd pay to avoid risk.

\begin{formulabox}
\textbf{Risk Premium:}
\[
\boxed{RP = EV(L) - CE}
\]

This is the maximum you'd pay to eliminate uncertainty!
\end{formulabox}

\begin{examplebox}
\textbf{Continuing previous example:}

$EV(L) = \$50$, $CE = \$25$

\textbf{Risk premium:}
\[
RP = 50 - 25 = \boxed{\$25}
\]

You'd pay up to \$25 to avoid the gamble and get the expected value for sure.
\end{examplebox}

\subsubsection{Relationship Between CE, RP, and EV}

\begin{tcolorbox}[colback=blue!5,colframe=blue!70]
\textbf{For risk-averse individuals:}
\[
CE < EV(L) \quad \text{and} \quad RP > 0
\]

\textbf{For risk-neutral individuals:}
\[
CE = EV(L) \quad \text{and} \quad RP = 0
\]

\textbf{For risk-loving individuals:}
\[
CE > EV(L) \quad \text{and} \quad RP < 0
\]
\end{tcolorbox}

\subsection{Common Utility Functions}

\subsubsection{1. Square Root: $u(x) = \sqrt{x}$}

\textbf{Properties:}
\begin{itemize}
  \item $u'(x) = \frac{1}{2\sqrt{x}} > 0$ (increasing)
  \item $u''(x) = -\frac{1}{4x^{3/2}} < 0$ (concave, risk averse)
  \item Easy to work with!
\end{itemize}

\subsubsection{2. Logarithmic: $u(x) = \ln(x)$}

\textbf{Properties:}
\begin{itemize}
  \item $u'(x) = \frac{1}{x} > 0$ (increasing)
  \item $u''(x) = -\frac{1}{x^2} < 0$ (concave, risk averse)
  \item Very risk averse for small $x$
\end{itemize}

\subsubsection{3. CARA: $u(x) = -e^{-ax}$ where $a > 0$}

\textbf{Properties:}
\begin{itemize}
  \item Constant Absolute Risk Aversion
  \item $u'(x) = ae^{-ax} > 0$ (increasing)
  \item $u''(x) = -a^2e^{-ax} < 0$ (concave, risk averse)
  \item Risk aversion doesn't depend on wealth level!
\end{itemize}

\subsubsection{4. CRRA: $u(x) = \frac{x^{1-\gamma}}{1-\gamma}$ where $\gamma > 0$}

\textbf{Properties:}
\begin{itemize}
  \item Constant Relative Risk Aversion
  \item When $\gamma = 1$: $u(x) = \ln(x)$
  \item When $\gamma = 0$: $u(x) = x$ (risk neutral)
  \item Commonly used in finance!
\end{itemize}

\subsection{Insurance}

\begin{criticalbox}
\textbf{STANDARD EXAM QUESTION:} Insurance problems test your understanding of risk aversion!
\end{criticalbox}

\subsubsection{The Setup}

You have initial wealth $W$ and face potential loss $L$ with probability $p$.

\textbf{Without insurance:}
\begin{itemize}
  \item Good state (no loss): wealth = $W$ with probability $1-p$
  \item Bad state (loss): wealth = $W - L$ with probability $p$
\end{itemize}

\textbf{Expected utility without insurance:}
\[
EU_{\text{no ins}} = (1-p) \cdot u(W) + p \cdot u(W-L)
\]

\textbf{With insurance:} Pay premium $\pi$ to get coverage $K$ if loss occurs
\begin{itemize}
  \item Good state: wealth = $W - \pi$ with probability $1-p$
  \item Bad state: wealth = $W - \pi - L + K$ with probability $p$
\end{itemize}

\textbf{Expected utility with insurance:}
\[
EU_{\text{ins}} = (1-p) \cdot u(W-\pi) + p \cdot u(W-\pi-L+K)
\]

\subsubsection{Actuarially Fair Insurance}

Insurance is \textbf{actuarially fair} if the premium equals expected payout:
\[
\boxed{\pi = p \cdot K}
\]

\textbf{Big Result:} Risk-averse people will buy full insurance ($K = L$) if it's actuarially fair!

\begin{examplebox}
\textbf{Problem:} You have $W = \$100$ and face 25\% chance of losing $L = \$60$. Your utility is $u(x) = \sqrt{x}$. An insurance company offers to cover the loss for premium $\pi = \$15$.

\textbf{(a)} Is this insurance actuarially fair?

\textbf{Solution (a):}
\[
\text{Fair premium} = p \cdot L = 0.25 \times 60 = \$15
\]

Yes! $\pi = \$15$ equals expected payout.

\textbf{(b)} Should you buy it?

\textbf{Solution (b):}

\textbf{Without insurance:}
\[
EU_{\text{no}} = 0.75 \times \sqrt{100} + 0.25 \times \sqrt{40} = 0.75 \times 10 + 0.25 \times 6.32 \approx 9.08
\]

\textbf{With insurance:}
\[
EU_{\text{ins}} = 1.0 \times \sqrt{100 - 15} = \sqrt{85} \approx 9.22
\]

Since $9.22 > 9.08$, \textbf{YES, buy insurance!}

You end up with \$85 for sure, which you prefer to the risky lottery.
\end{examplebox}

\subsection{Portfolio Choice}

\subsubsection{The Setup}

You have wealth $W$ to allocate between:
\begin{itemize}
  \item \textbf{Safe asset:} Returns $r_f$ (certain)
  \item \textbf{Risky asset:} Returns $r_h$ with probability $p$, $r_l$ with probability $1-p$ where $r_h > r_f > r_l$
\end{itemize}

Let $\alpha$ = fraction invested in risky asset.

\textbf{Final wealth:}
\begin{itemize}
  \item Good state: $W_h = W[(1-\alpha)(1+r_f) + \alpha(1+r_h)]$
  \item Bad state: $W_l = W[(1-\alpha)(1+r_f) + \alpha(1+r_l)]$
\end{itemize}

\textbf{Expected utility:}
\[
EU(\alpha) = p \cdot u(W_h) + (1-p) \cdot u(W_l)
\]

\textbf{Goal:} Choose $\alpha^*$ to maximize $EU(\alpha)$!

\subsubsection{First-Order Condition}

\begin{formulabox}
\textbf{Optimal Portfolio FOC:}
\[
\boxed{\frac{dEU}{d\alpha} = p \cdot u'(W_h) \cdot \frac{dW_h}{d\alpha} + (1-p) \cdot u'(W_l) \cdot \frac{dW_l}{d\alpha} = 0}
\]

Where:
\begin{align*}
\frac{dW_h}{d\alpha} &= W(r_h - r_f) \\
\frac{dW_l}{d\alpha} &= W(r_l - r_f)
\end{align*}
\end{formulabox}

\begin{examplebox}
\textbf{Problem:} You have $W = \$100$ and $u(x) = \ln(x)$. Safe asset returns 0\%, risky asset returns +20\% or -10\% with equal probability. Find optimal $\alpha$.

\textbf{Solution:}

$r_f = 0$, $r_h = 0.2$, $r_l = -0.1$, $p = 0.5$

\textbf{Final wealth:}
\begin{align*}
W_h &= 100[(1-\alpha)(1) + \alpha(1.2)] = 100(1 + 0.2\alpha) \\
W_l &= 100[(1-\alpha)(1) + \alpha(0.9)] = 100(1 - 0.1\alpha)
\end{align*}

\textbf{FOC:}
\[
0.5 \cdot \frac{1}{W_h} \cdot 20 + 0.5 \cdot \frac{1}{W_l} \cdot (-10) = 0
\]

\[
\frac{10}{100(1+0.2\alpha)} = \frac{5}{100(1-0.1\alpha)}
\]

\[
\frac{10}{1+0.2\alpha} = \frac{5}{1-0.1\alpha}
\]

Cross-multiply:
\[
10(1-0.1\alpha) = 5(1+0.2\alpha)
\]
\[
10 - \alpha = 5 + \alpha
\]
\[
5 = 2\alpha \implies \boxed{\alpha^* = 2.5 = 250\%}
\]

Wait, this is >100\%! This means \textbf{leverage} - borrow money to invest more in risky asset!
\end{examplebox}

\warningbox{
\textbf{Watch out for corner solutions!}

If FOC gives $\alpha > 1$, check if borrowing is allowed.

If FOC gives $\alpha < 0$, check if short-selling is allowed.

If neither allowed, optimal is $\alpha^* = 0$ or $\alpha^* = 1$ (corner solution).
}

\newpage

\section{PRODUCER THEORY - Firms and Costs}

\begin{criticalbox}
\textbf{Why This Matters:} Foundation for understanding markets!

Master cost functions and you'll breeze through monopoly and competition problems.
\end{criticalbox}

\subsection{Production Functions}

\subsubsection{Definition}

A \textbf{production function} shows how inputs (labor $L$, capital $K$) are transformed into output $q$.

\begin{formulabox}
\textbf{Production Function:}
\[
\boxed{q = f(L, K)}
\]

\textbf{Common examples:}
\begin{itemize}
  \item Linear: $q = aL + bK$
  \item Cobb-Douglas: $q = AL^\alpha K^\beta$
  \item Square root: $q = \sqrt{L}$ or $q = a + \sqrt{L}$
\end{itemize}
\end{formulabox}

\subsection{Cost Functions}

\subsubsection{From Production to Cost}

\textbf{The process:}

\textbf{Step 1:} Start with production function $q = f(L, K)$

\textbf{Step 2:} Invert to express input in terms of output
\begin{itemize}
  \item Example: $q = \sqrt{L} \implies L = q^2$
\end{itemize}

\textbf{Step 3:} Multiply by input prices
\begin{itemize}
  \item If labor costs $w$ per unit: $C(q) = w \cdot L(q)$
  \item If capital costs $r$ per unit: add $r \cdot K(q)$
\end{itemize}

\begin{examplebox}
\textbf{Problem:} Firm has $q = 2\sqrt{L}$ where $w = \$9$ per unit. Find $C(q)$.

\textbf{Solution:}

\textbf{Step 1:} Invert production function
\[
q = 2\sqrt{L} \implies \frac{q}{2} = \sqrt{L} \implies L = \frac{q^2}{4}
\]

\textbf{Step 2:} Multiply by wage
\[
C(q) = w \cdot L = 9 \cdot \frac{q^2}{4} = \boxed{\frac{9q^2}{4}}
\]
\end{examplebox}

\subsubsection{Cost Concepts}

\begin{formulabox}
\textbf{Total Cost:}
\[
C(q) = FC + VC(q)
\]
where $FC$ = fixed cost, $VC(q)$ = variable cost

\textbf{Marginal Cost (MC):}
\[
\boxed{MC(q) = \frac{dC}{dq} = C'(q)}
\]
Cost of producing one more unit.

\textbf{Average Cost (AC):}
\[
\boxed{AC(q) = \frac{C(q)}{q}}
\]
Cost per unit.

\textbf{Average Variable Cost (AVC):}
\[
\boxed{AVC(q) = \frac{VC(q)}{q}}
\]
\end{formulabox}

\subsubsection{Key Relationships}

\begin{tcolorbox}[colback=blue!5,colframe=blue!70]
\textbf{Important facts:}
\begin{itemize}
  \item MC intersects AC at minimum of AC
  \item If MC < AC, then AC is decreasing
  \item If MC > AC, then AC is increasing
  \item Firms shut down if $P < \min AVC$ (short run)
\end{itemize}
\end{tcolorbox}

\newpage

\section{PERFECT COMPETITION vs. MONOPOLY}

\begin{criticalbox}
\textbf{KEY DISTINCTION:} Perfect competition → price takers. Monopoly → price setters.

This difference drives all results!
\end{criticalbox}

\subsection{Perfect Competition}

\subsubsection{Characteristics}

\begin{itemize}
  \item Many small firms
  \item Homogeneous product (all firms sell identical goods)
  \item Free entry and exit
  \item Perfect information
\end{itemize}

\textbf{Result:} Firms are \textbf{price takers} - they take market price $P$ as given.

\subsubsection{Firm's Problem}

\begin{formulabox}
\textbf{Competitive Firm's Profit:}
\[
\boxed{\Pi(q) = P \cdot q - C(q)}
\]

\textbf{FOC (profit maximization):}
\[
\boxed{P = MC(q)}
\]

Firms produce where price equals marginal cost!
\end{formulabox}

\subsubsection{Supply Curve}

A competitive firm's supply curve is its MC curve above minimum AVC.

\textbf{Shutdown rule:} If $P < \min AVC$, produce $q = 0$ (shut down).

\subsubsection{Long-Run Equilibrium}

In long run, free entry/exit drives profits to zero:

\begin{formulabox}
\textbf{Long-Run Equilibrium Conditions:}
\begin{enumerate}
  \item $P = MC$ (profit maximization)
  \item $P = AC$ (zero profit condition from free entry)
  \item Therefore: $\boxed{P = MC = \min AC}$
\end{enumerate}

Firms produce at minimum efficient scale!
\end{formulabox}

\begin{examplebox}
\textbf{Problem:} A competitive firm has $C(q) = q^2 + 4q + 4$. Market price is $P = \$12$. Find optimal output.

\textbf{Solution:}

\textbf{Step 1:} Find MC
\[
MC(q) = C'(q) = 2q + 4
\]

\textbf{Step 2:} Set $P = MC$
\[
12 = 2q + 4 \implies 8 = 2q \implies \boxed{q^* = 4}
\]

\textbf{Check profit:}
\[
\Pi = P \cdot q - C(q) = 12 \times 4 - (16 + 16 + 4) = 48 - 36 = \$12 > 0
\]

Firm makes positive profit! (But in long run, entry would drive profit to zero)
\end{examplebox}

\subsection{Monopoly}

\subsubsection{Characteristics}

\begin{itemize}
  \item Single firm (only seller)
  \item No close substitutes
  \item Barriers to entry (prevent competition)
\end{itemize}

\textbf{Result:} Firm is \textbf{price setter} - it faces entire market demand curve.

\subsubsection{Demand and Revenue}

The monopolist faces market demand $Q^D(P)$ or inverse demand $P^D(Q)$.

\begin{formulabox}
\textbf{Monopolist's Revenue:}
\[
R(Q) = P^D(Q) \cdot Q
\]

\textbf{Marginal Revenue:}
\[
\boxed{MR(Q) = \frac{dR}{dQ} = P^D(Q) + Q \cdot \frac{dP^D}{dQ}}
\]

Since $\frac{dP^D}{dQ} < 0$ (downward sloping demand), we have:
\[
\boxed{MR(Q) < P^D(Q)}
\]

\textbf{Marginal revenue is less than price!}
\end{formulabox}

\subsubsection{Why MR < P?}

To sell one more unit, monopolist must lower price. This has two effects:
\begin{enumerate}
  \item \textbf{Gain:} Revenue from extra unit sold at price $P$ ($+P$)
  \item \textbf{Loss:} Lower price on all previous units ($-Q \cdot \Delta P$)
\end{enumerate}

Net effect: $MR = P + Q \cdot \frac{dP}{dQ} < P$

\subsubsection{Monopoly Profit Maximization}

\begin{formulabox}[colback=yellow!20,colframe=orange!80]
\textbf{Monopolist's Problem:}
\[
\max_Q \Pi(Q) = R(Q) - C(Q) = P^D(Q) \cdot Q - C(Q)
\]

\textbf{FOC:}
\[
\boxed{MR(Q) = MC(Q)}
\]

\textbf{Steps to solve:}
\begin{enumerate}
  \item Find $MR(Q)$ from demand
  \item Find $MC(Q)$ from cost function
  \item Set $MR = MC$ and solve for $Q^*$
  \item Find price: $P^* = P^D(Q^*)$
\end{enumerate}
\end{formulabox}

\subsubsection{MR Formula for Linear Demand}

\begin{tcolorbox}[colback=yellow!10,colframe=orange]
\textbf{SHORTCUT:} If inverse demand is linear $P = a - bQ$, then:
\[
\boxed{MR = a - 2bQ}
\]

MR has same intercept but \textbf{twice the slope}!
\end{tcolorbox}

\begin{examplebox}
\textbf{Problem:} Monopolist has $C(Q) = Q^2 + 10$ and faces demand $Q^D(P) = 50 - P$. Find optimal quantity and price.

\textbf{Solution:}

\textbf{Step 1:} Convert to inverse demand
\[
Q = 50 - P \implies P^D(Q) = 50 - Q
\]

\textbf{Step 2:} Find MR
\[
MR(Q) = 50 - 2Q \quad \text{(double the slope!)}
\]

\textbf{Step 3:} Find MC
\[
MC(Q) = C'(Q) = 2Q
\]

\textbf{Step 4:} Set $MR = MC$
\[
50 - 2Q = 2Q \implies 50 = 4Q \implies \boxed{Q^* = 12.5}
\]

\textbf{Step 5:} Find price
\[
P^* = P^D(12.5) = 50 - 12.5 = \boxed{P^* = 37.5}
\]

\textbf{Profit:}
\[
\Pi = 37.5 \times 12.5 - (12.5^2 + 10) = 468.75 - 166.25 = \boxed{\$302.50}
\]
\end{examplebox}

\subsection{Comparing Perfect Competition and Monopoly}

\begin{table}[h]
\centering
\begin{tabular}{lcc}
\toprule
\textbf{Feature} & \textbf{Perfect Competition} & \textbf{Monopoly} \\
\midrule
Number of firms & Many & One \\
Profit max condition & $P = MC$ & $MR = MC$ \\
Price vs. MC & $P = MC$ & $P > MC$ \\
Output level & Higher & Lower \\
Price level & Lower & Higher \\
Profit (long run) & Zero (free entry) & Positive (barriers) \\
Efficiency & Efficient & Inefficient (DWL) \\
\bottomrule
\end{tabular}
\end{table}

\textbf{Key insight:} Monopoly creates deadweight loss by restricting output to raise price!

\newpage

\section{GENERAL EQUILIBRIUM AND WELFARE}

\begin{criticalbox}
\textbf{BIG PICTURE:} How do markets interact? When are outcomes efficient?
\end{criticalbox}

\subsection{Partial vs. General Equilibrium}

\textbf{Partial equilibrium:} Analyze one market in isolation (what we've done so far)

\textbf{General equilibrium:} Analyze all markets simultaneously, accounting for interactions

\subsection{Exchange Economy - The Edgeworth Box}

\subsubsection{Setup}

Two consumers (A and B), two goods (X and Y).

\textbf{Endowments:}
\begin{itemize}
  \item Consumer A starts with $(\omega_A^X, \omega_A^Y)$
  \item Consumer B starts with $(\omega_B^X, \omega_B^Y)$
\end{itemize}

\textbf{Total resources:} $\bar{X} = \omega_A^X + \omega_B^X$, $\bar{Y} = \omega_A^Y + \omega_B^Y$

\textbf{Question:} Can they trade to make both better off?

\subsubsection{Pareto Efficiency}

An allocation is \textbf{Pareto efficient} if you can't make someone better off without making someone else worse off.

\begin{formulabox}
\textbf{Pareto Efficiency Condition:}
\[
\boxed{MRS_A = MRS_B}
\]

Marginal rates of substitution must be equal!

If $MRS_A \neq MRS_B$, there are gains from trade.
\end{formulabox}

\subsubsection{First Welfare Theorem}

\begin{tcolorbox}[colback=green!5,colframe=green!70]
\textbf{FIRST WELFARE THEOREM:}

If markets are competitive and complete, the equilibrium is Pareto efficient.

\textbf{Translation:} Free markets lead to efficient outcomes (under ideal conditions)!
\end{tcolorbox}

\subsubsection{Second Welfare Theorem}

\begin{tcolorbox}[colback=green!5,colframe=green!70]
\textbf{SECOND WELFARE THEOREM:}

Any Pareto efficient allocation can be achieved as a competitive equilibrium with appropriate redistribution.

\textbf{Translation:} Can achieve any efficient outcome by redistributing endowments, then letting markets work!
\end{tcolorbox}

\subsection{Consumer and Producer Surplus}

\subsubsection{Consumer Surplus (CS)}

The difference between what consumers are willing to pay and what they actually pay.

\textbf{Graphically:} Area under demand curve, above price.

\begin{formulabox}
\textbf{Consumer Surplus:}
\[
\boxed{CS = \int_0^Q P^D(q) \, dq - P \cdot Q}
\]
\end{formulabox}

\subsubsection{Producer Surplus (PS)}

The difference between what producers receive and their marginal cost.

\textbf{Graphically:} Area above supply curve (MC), below price.

\begin{formulabox}
\textbf{Producer Surplus:}
\[
\boxed{PS = P \cdot Q - \int_0^Q MC(q) \, dq}
\]
\end{formulabox}

\subsubsection{Total Surplus and Deadweight Loss}

\begin{formulabox}
\textbf{Total Surplus:}
\[
\boxed{TS = CS + PS}
\]

This measures total welfare in the market.

\textbf{Deadweight Loss (DWL):}
\[
\boxed{DWL = TS_{\text{efficient}} - TS_{\text{actual}}}
\]

Lost welfare from inefficiency (e.g., monopoly).
\end{formulabox}

\textbf{Perfect competition:} $DWL = 0$ (efficient)

\textbf{Monopoly:} $DWL > 0$ (inefficient - monopolist restricts output)

\newpage

\section{PATTERN RECOGNITION GUIDE}

\begin{strategybox}
\textbf{How to identify problem types and choose the right approach}
\end{strategybox}

\subsection{Uncertainty Problems}

\begin{table}[h]
\centering
\begin{tabular}{p{4cm}p{4cm}p{6cm}}
\toprule
\textbf{Key Words} & \textbf{What to Do} & \textbf{Formula to Use} \\
\midrule
``Expected value'' & Calculate weighted average & $EV = \sum p_i x_i$ \\
\midrule
``Expected utility'' & Calculate $\sum p_i u(x_i)$ & $EU = \sum p_i u(x_i)$ \\
\midrule
``Certainty equivalent'' & Solve $u(CE) = EU$ & $u(CE) = EU(L)$ \\
\midrule
``Risk premium'' & Find $EV - CE$ & $RP = EV - CE$ \\
\midrule
``Risk averse'' & Check concavity & $u''(x) < 0$ \\
\midrule
``Insurance'' & Compare EU with/without & Set up lottery, calculate both \\
\midrule
``Portfolio'' & Maximize $EU(\alpha)$ & Take FOC: $\frac{dEU}{d\alpha} = 0$ \\
\bottomrule
\end{tabular}
\end{table}

\subsection{Producer/Cost Problems}

\begin{table}[h]
\centering
\begin{tabular}{p{4cm}p{4cm}p{6cm}}
\toprule
\textbf{Key Words} & \textbf{What to Do} & \textbf{Formula to Use} \\
\midrule
``Production function'' & Invert to get input(output) & $q = f(L) \implies L(q)$ \\
\midrule
``Cost function'' & Multiply input by price & $C(q) = w \cdot L(q)$ \\
\midrule
``Marginal cost'' & Take derivative of $C(q)$ & $MC = C'(q)$ \\
\midrule
``Average cost'' & Divide cost by quantity & $AC = C(q)/q$ \\
\bottomrule
\end{tabular}
\end{table}

\subsection{Market Structure Problems}

\begin{table}[h]
\centering
\begin{tabular}{p{4cm}p{4cm}p{6cm}}
\toprule
\textbf{Key Words} & \textbf{Type} & \textbf{Profit Max Rule} \\
\midrule
``Competitive firm'', ``price taker'' & Perfect Competition & $P = MC$ \\
\midrule
``Monopolist'', ``only firm'', ``faces demand'' & Monopoly & $MR = MC$ \\
\midrule
``Long-run equilibrium'', ``free entry'' & Long-run Competition & $P = MC = \min AC$ \\
\bottomrule
\end{tabular}
\end{table}

\subsection{Decision Tree for Solving Problems}

\begin{enumerate}
  \item \textbf{Is there uncertainty?}
    \begin{itemize}
      \item YES → Go to Section 1 (Uncertainty)
      \item NO → Continue to step 2
    \end{itemize}

  \item \textbf{Is it about a firm?}
    \begin{itemize}
      \item YES → Continue to step 3
      \item NO → Check if it's general equilibrium (Section 4)
    \end{itemize}

  \item \textbf{What market structure?}
    \begin{itemize}
      \item Many firms, price taker → Perfect competition (Section 3.1)
      \item One firm, faces demand → Monopoly (Section 3.2)
    \end{itemize}

  \item \textbf{What's the question?}
    \begin{itemize}
      \item Find cost function → Invert production function (Section 2.2)
      \item Find optimal quantity → Set MR = MC (Section 3)
      \item Find price → Use demand curve (monopoly) or given (competition)
    \end{itemize}
\end{enumerate}

\newpage

\section{FORMULA SHEET - EVERYTHING YOU NEED}

\subsection{Uncertainty Formulas}

\begin{tcolorbox}[colback=formulacolor!5,colframe=formulacolor,title={\textbf{Expected Value \& Utility}}]
\textbf{Expected Value:}
\[
EV(L) = \sum_{i=1}^n p_i x_i
\]

\textbf{Expected Utility:}
\[
EU(L) = \sum_{i=1}^n p_i \cdot u(x_i)
\]

\textbf{Certainty Equivalent:}
\[
u(CE) = EU(L)
\]

\textbf{Risk Premium:}
\[
RP = EV(L) - CE
\]

\textbf{Risk Aversion:}
\begin{itemize}
  \item Risk averse: $u''(x) < 0$ (concave), $CE < EV$, $RP > 0$
  \item Risk neutral: $u''(x) = 0$ (linear), $CE = EV$, $RP = 0$
  \item Risk loving: $u''(x) > 0$ (convex), $CE > EV$, $RP < 0$
\end{itemize}
\end{tcolorbox}

\subsection{Common Utility Functions}

\begin{tcolorbox}[colback=formulacolor!5,colframe=formulacolor,title={\textbf{Utility Functions}}]
\textbf{Square Root:} $u(x) = \sqrt{x}$
\begin{itemize}
  \item $u'(x) = \frac{1}{2\sqrt{x}}$, $u''(x) = -\frac{1}{4x^{3/2}} < 0$ (risk averse)
\end{itemize}

\textbf{Logarithmic:} $u(x) = \ln(x)$
\begin{itemize}
  \item $u'(x) = \frac{1}{x}$, $u''(x) = -\frac{1}{x^2} < 0$ (risk averse)
\end{itemize}

\textbf{CARA:} $u(x) = -e^{-ax}$ where $a > 0$
\begin{itemize}
  \item Constant absolute risk aversion
\end{itemize}

\textbf{CRRA:} $u(x) = \frac{x^{1-\gamma}}{1-\gamma}$ where $\gamma > 0$
\begin{itemize}
  \item Constant relative risk aversion
  \item When $\gamma = 1$: $u(x) = \ln(x)$
\end{itemize}
\end{tcolorbox}

\subsection{Cost Functions}

\begin{tcolorbox}[colback=formulacolor!5,colframe=formulacolor,title={\textbf{Costs}}]
\textbf{Total Cost:}
\[
C(q) = FC + VC(q)
\]

\textbf{Marginal Cost:}
\[
MC(q) = \frac{dC}{dq} = C'(q)
\]

\textbf{Average Cost:}
\[
AC(q) = \frac{C(q)}{q}
\]

\textbf{Average Variable Cost:}
\[
AVC(q) = \frac{VC(q)}{q}
\]
\end{tcolorbox}

\subsection{Perfect Competition}

\begin{tcolorbox}[colback=formulacolor!5,colframe=formulacolor,title={\textbf{Competitive Firms}}]
\textbf{Profit:}
\[
\Pi(q) = P \cdot q - C(q)
\]

\textbf{Profit Maximization:}
\[
P = MC(q)
\]

\textbf{Shutdown Rule:}
\[
\text{Produce if } P \geq \min AVC, \text{ else } q = 0
\]

\textbf{Long-Run Equilibrium:}
\[
P = MC = \min AC \quad \text{(zero profit)}
\]
\end{tcolorbox}

\subsection{Monopoly}

\begin{tcolorbox}[colback=yellow!20,colframe=orange!80,title={\textbf{Monopoly - MOST IMPORTANT}}]
\textbf{Revenue:}
\[
R(Q) = P^D(Q) \cdot Q
\]

\textbf{Marginal Revenue:}
\[
MR(Q) = P^D(Q) + Q \cdot \frac{dP^D}{dQ}
\]

\textbf{For linear demand} $P = a - bQ$:
\[
\boxed{MR = a - 2bQ} \quad \text{(twice the slope!)}
\]

\textbf{Profit Maximization:}
\[
MR(Q) = MC(Q)
\]

\textbf{Then find price:}
\[
P^* = P^D(Q^*)
\]

\textbf{Key fact:}
\[
MR < P \quad \text{(MR is always less than price)}
\]
\end{tcolorbox}

\subsection{Welfare}

\begin{tcolorbox}[colback=formulacolor!5,colframe=formulacolor,title={\textbf{Consumer \& Producer Surplus}}]
\textbf{Consumer Surplus:}
\[
CS = \int_0^Q P^D(q) \, dq - P \cdot Q
\]
(Area under demand, above price)

\textbf{Producer Surplus:}
\[
PS = P \cdot Q - \int_0^Q MC(q) \, dq
\]
(Area above MC, below price)

\textbf{Total Surplus:}
\[
TS = CS + PS
\]

\textbf{Deadweight Loss:}
\[
DWL = TS_{\text{efficient}} - TS_{\text{actual}}
\]
\end{tcolorbox}

\subsection{General Equilibrium}

\begin{tcolorbox}[colback=formulacolor!5,colframe=formulacolor,title={\textbf{Pareto Efficiency}}]
\textbf{Pareto Efficiency Condition:}
\[
MRS_A = MRS_B
\]

\textbf{First Welfare Theorem:} Competitive equilibrium is Pareto efficient

\textbf{Second Welfare Theorem:} Any Pareto efficient allocation can be achieved as competitive equilibrium with redistribution
\end{tcolorbox}

\newpage

\section{COMMON MISTAKES TO AVOID}

\begin{warningbox}
\textbf{Mistake 1: Confusing EV and EU}

Don't use $EV = \sum p_i x_i$ when asked for expected utility!

✓ Correct: $EU = \sum p_i u(x_i)$
\end{warningbox}

\begin{warningbox}
\textbf{Mistake 2: Forgetting to invert production function}

Can't just replace $L$ with $q$ in production function!

✗ Wrong: $q = \sqrt{L} \implies C(q) = w\sqrt{q}$

✓ Correct: $q = \sqrt{L} \implies L = q^2 \implies C(q) = wq^2$
\end{warningbox}

\begin{warningbox}
\textbf{Mistake 3: Using $P = MC$ for monopoly}

Monopolists set $MR = MC$, not $P = MC$!

✗ Wrong: $P = MC$ (this is perfect competition)

✓ Correct: $MR = MC$ (for monopoly)
\end{warningbox}

\begin{warningbox}
\textbf{Mistake 4: Wrong MR for linear demand}

If $P = a - bQ$, then $MR = a - 2bQ$ (not $a - bQ$!)

Remember: MR has twice the slope of demand!
\end{warningbox}

\begin{warningbox}
\textbf{Mistake 5: Forgetting demand vs. inverse demand}

$Q^D(P)$ tells you quantity at given price.

$P^D(Q)$ tells you price at given quantity.

For monopoly revenue: $R(Q) = P^D(Q) \cdot Q$ (use inverse demand!)
\end{warningbox}

\begin{warningbox}
\textbf{Mistake 6: Not checking for corner solutions}

If portfolio FOC gives $\alpha > 1$ or $\alpha < 0$, check constraints!

May need corner solution: $\alpha^* = 0$ or $\alpha^* = 1$
\end{warningbox}

\newpage

\section*{Final Checklist - Are You Ready?}

\begin{tcolorbox}[colback=green!10,colframe=green!70,arc=3mm]
\textbf{Before the exam, make sure you can:}

\subsection*{Uncertainty}
\begin{itemize}
  \item[$\square$] Calculate expected value and expected utility
  \item[$\square$] Find certainty equivalent by solving $u(CE) = EU$
  \item[$\square$] Calculate risk premium: $RP = EV - CE$
  \item[$\square$] Determine if utility function is risk averse (check $u''(x) < 0$)
  \item[$\square$] Solve insurance problems (compare EU with/without insurance)
  \item[$\square$] Solve portfolio choice problems (maximize $EU(\alpha)$ with FOC)
\end{itemize}

\subsection*{Producer Theory}
\begin{itemize}
  \item[$\square$] Invert production function to get input as function of output
  \item[$\square$] Derive cost function: $C(q) = w \cdot L(q)$
  \item[$\square$] Calculate marginal cost: $MC = C'(q)$
  \item[$\square$] Calculate average cost: $AC = C(q)/q$
\end{itemize}

\subsection*{Perfect Competition}
\begin{itemize}
  \item[$\square$] Set $P = MC$ to find optimal quantity
  \item[$\square$] Apply shutdown rule: produce if $P \geq \min AVC$
  \item[$\square$] Understand long-run equilibrium: $P = MC = \min AC$
\end{itemize}

\subsection*{Monopoly}
\begin{itemize}
  \item[$\square$] Convert demand to inverse demand
  \item[$\square$] Calculate MR (for linear demand: twice the slope!)
  \item[$\square$] Set $MR = MC$ to find optimal $Q$
  \item[$\square$] Find price using demand curve: $P^* = P^D(Q^*)$
  \item[$\square$] Remember: $MR < P$ always for monopolist!
\end{itemize}

\subsection*{General Equilibrium \& Welfare}
\begin{itemize}
  \item[$\square$] Understand Pareto efficiency: $MRS_A = MRS_B$
  \item[$\square$] Know First Welfare Theorem (competition → efficiency)
  \item[$\square$] Calculate consumer and producer surplus
  \item[$\square$] Understand deadweight loss from monopoly
\end{itemize}
\end{tcolorbox}

\vspace{1cm}

\begin{center}
\begin{tcolorbox}[colback=yellow!20,colframe=orange!80,width=0.9\textwidth,arc=5mm]
\centering
\LARGE\textbf{You've got this!}

\large Master the formulas, practice the problems, and trust your preparation.

\textbf{Good luck on the final exam!}
\end{tcolorbox}
\end{center}

\end{document}
